\section{Equilibrium states and vanishing replica cross terms}
\label{sec:equilibrium}

The spatial estimate has a second use besides comparing block averages:
it removes certain couplings between independent replicas at a pressure
tangent.  We give the variational framework first, including the precise
order of limits needed for this deduction.  The underlying lattice is
now $\Z^3$, with the spin-$S$ matrix algebra at each site.  A state is a
consistent family of density matrices on finite sets, or equivalently
a positive normalized functional on the norm closure of the local
matrix algebras.  Local convergence means convergence on every local
observable.  Write $\mathcal S_{\rm TI}$ for the compact convex set of
translation-invariant states, and $\tau_x$ for translation by $x$.

\subsection{The variational principle and its tangents}

For a density matrix $\rho$, let
$\mathsf S(\rho)=-\Tr(\rho\log\rho)$.  If $Q_n=\{0,\ldots,n-1\}^3$
and $\omega\in\mathcal S_{\rm TI}$, its mean entropy and unperturbed
interaction expectation per site are
\[
 \mathfrak s(\omega)=\lim_{n\to\infty}\frac{\mathsf S(\omega_{Q_n})}{|Q_n|},
 \qquad
 e_0(\omega)=\beta\sum_{r=1}^3
                     \omega(\mathbf S_0\cdot\mathbf S_{e_r}).
\]
For a self-adjoint local observable $a$, let
$A_L=\sum_{x\in\Lambda_L}\tau_x(a)$, interpreted periodically for $L$
larger than its support, and set
\[
 P_a(t)=\lim_{L\to\infty}\frac1V
                   \log\Tr e^{H_0+tA_L},\qquad t\in\R.
\]
The limit uses even $L$, as elsewhere.  Finite range and boundedness
make free and periodic boundary conditions equivalent for this limit:
the difference of their exponents has norm $O(L^2)$.

The following standard variational facts are recalled with proofs to
specify the state space on which the later tangent is taken; see
\cite[Theorems~2--3]{LanfordRobinson1968Entropy} and
\cite[Theorem~3]{LanfordRobinson1968Spin}.

\begin{proposition}\label{prop:variational-framework}
The mean entropy exists on $\mathcal S_{\rm TI}$, is upper
semicontinuous and affine, and
\begin{equation}\label{eq:variational-pressure}
 P_a(t)=\max_{\omega\in\mathcal S_{\rm TI}}
       \{\mathfrak s(\omega)+e_0(\omega)+t\omega(a)\}.
\end{equation}
In particular, writing $P_0=\beta p_{S,\beta}(0)$, the set
\[
 \mathcal E_\beta=
 \{\omega\in\mathcal S_{\rm TI}:
            \mathfrak s(\omega)+e_0(\omega)=P_0\}
\]
is a nonempty compact face of $\mathcal S_{\rm TI}$.  Every local limit
of periodic zero-field Gibbs states belongs to $\mathcal E_\beta$,
and
\begin{equation}\label{eq:pressure-tangent}
 P_a'(0+)=\max_{\omega\in\mathcal E_\beta}\omega(a).
\end{equation}
The extreme points of $\mathcal E_\beta$ are precisely its states that
are extreme in $\mathcal S_{\rm TI}$; we call these states ergodic.
\end{proposition}

\begin{proof}
Tile a large cube by translates of $Q_n$ and a remainder of vanishing
relative volume.  Entropy subadditivity and the bound
$\mathsf S(\rho_B)\le |B|\log(2S+1)$ show that
\[
 \mathfrak s(\omega)
       =\inf_{n\ge1}\frac{\mathsf S(\omega_{Q_n})}{|Q_n|}.
\]
This proves existence and upper semicontinuity, since each finite-volume
entropy is continuous.  For a mixture $\rho=\lambda\rho_1+(1-\lambda)\rho_2$,
the entropy lies between
$\lambda\mathsf S(\rho_1)+(1-\lambda)\mathsf S(\rho_2)$ and this quantity
plus the binary entropy of $\lambda$.  Division by volume proves
affinity of $\mathfrak s$.

Apply the finite-dimensional Gibbs variational principle to the
restriction of any $\omega\in\mathcal S_{\rm TI}$ to a large cube.
Boundary terms are $o(V)$ in norm, so
$P_a(t)\ge\mathfrak s(\omega)+e_0(\omega)+t\omega(a)$.
Conversely, choose a locally convergent subsequence of periodic Gibbs
states for $H_0+tA_L$, with limit $\omega$.  The limit is translation
invariant, and local interaction expectations converge.  Tiling the
large cubes by fixed cubes $Q_n$ gives
\[
 \limsup_{L\to\infty}\frac{\mathsf S(\gamma_{L,t})}{V}
 \le \frac{\mathsf S(\omega_{Q_n})}{|Q_n|}
 \quad\hbox{for every }n.
\]
Letting $n\to\infty$ and using the exact finite-volume Gibbs identity
gives the reverse variational inequality.  At $t=0$ this argument also
proves the assertion about every periodic Gibbs local limit.  Compactness,
upper semicontinuity, and affinity show that $\mathcal E_\beta$ is a
nonempty compact face; the assertion about its extreme points follows
from the definition of a face.

For completeness, a maximizer $\omega_t$ in
\eqref{eq:variational-pressure}, with $t>0$, satisfies
\[
 \max_{\omega\in\mathcal E_\beta}\omega(a)
 \le\frac{P_a(t)-P_0}{t}\le\omega_t(a).
\]
As $t\downarrow0$, upper semicontinuity shows that every local limit
of $\omega_t$ belongs to $\mathcal E_\beta$: its unperturbed variational
value tends to $P_0$, since $|\omega_t(a)|\le\|a\|$.
Taking a convergent subsequence proves \eqref{eq:pressure-tangent}.
\end{proof}

In the field direction $Y$, the exponent-normalized pressure is
$\beta p_{S,\beta}(b/\beta)$.  Thus
\begin{equation}\label{eq:magnetization-max}
 m=\max_{\omega\in\mathcal E_\beta}\omega(Y_0).
\end{equation}
Rotation invariance gives the same maximum in every unit direction.
Consequently the vector
$\mathbf m(\omega)=(\omega(X_0),\omega(Y_0),\omega(Z_0))$ has length
at most $m$ for every $\omega\in\mathcal E_\beta$.  The face maximizing
$\omega(Y_0)$ has an extreme point, so an ergodic equilibrium state
attains $m$.  Section~\ref{sec:pressure} will prove that every ergodic
equilibrium state has this same magnetization magnitude.

The following mean-ergodic consequence of extremality is standard for
asymptotically commuting lattice translations; see
\cite[Theorem~3]{KastlerRobinson1966}.

\begin{lemma}[Ergodic averages]\label{lem:ergodic-variance}
If $\omega$ is an extreme translation-invariant state and $A$ is a
self-adjoint local observable, then
\[
 \lim_{n\to\infty}
 \omega\left[\left(
      \frac1{|Q_n|}\sum_{x\in Q_n}\tau_x(A)-\omega(A)\1
                   \right)^2\right]=0.
\]
\end{lemma}

\begin{proof}
After shifting and scaling, assume $0\le A\le\1$, and denote its
average by $\bar A_n$.  In the Hilbert space of the state, translations
are commuting unitaries.  The mean ergodic theorem makes
$\bar A_n\Omega$ converge to a vector $\eta$, where $\Omega$ represents
the state.  Hence $\lambda(C)=\lim_n\omega(C\bar A_n)$ exists for every
local $C$.

For local $C\ge0$, locality gives
$\|[C^{1/2},\bar A_n]\|\to0$.  Therefore
$\lambda(C)=\lim_n\omega(C^{1/2}\bar A_nC^{1/2})$, and
$0\le\lambda(C)\le\omega(C)$.  The functional extends continuously,
is translation invariant by the boundary-to-volume estimate for the
averages, and satisfies $\lambda(\1)=\omega(A)$.  Extremality of
$\omega$ now gives $\lambda=\omega(A)\omega$: otherwise its normalized
parts in $\omega=\lambda+(\omega-\lambda)$ would give a nontrivial
convex decomposition into translation-invariant states.  Testing on
the dense set of local vectors identifies $\eta=\omega(A)\Omega$,
which is the asserted variance convergence.
\end{proof}

\subsection{Products of zero-sum observables in distinct replicas}

Take $j$ copies of the spin system, with superscripts $[r]$ denoting
the replica.  The unperturbed exponent is
$\sum_{r=1}^jH_0^{[r]}$, and its pressure is $jP_0$.  Every
translation-invariant equilibrium state for this interaction has each
marginal in $\mathcal E_\beta$.  Indeed, entropy subadditivity between
replicas bounds its variational value by the sum of the marginal
variational values, each at most $P_0$; equality of the total forces
equality for each marginal.  Conversely, products of states in
$\mathcal E_\beta$ are replicated equilibrium states, because their
entropies and interaction expectations add.

For real coefficients $v_i$ on finitely many qubit slots, define
\[
 D_v^a=\sum_i v_i s_i^a,
 \qquad B_v^a=P_*D_v^aP_*
\]
as a physical local operator.  On a site's symmetric subspace,
$P_*s_{(x,k)}^aP_*=S_x^a/\ell$.  Thus
\begin{equation}\label{eq:compressed-zero-sum}
 B_v^a=\sum_x c_x S_x^a,
 \qquad c_x=\frac1\ell\sum_{k=1}^{\ell}v_{(x,k)},
 \qquad \sum_xc_x=0\quad\hbox{if }\sum_iv_i=0.
\end{equation}
The next lemma concerns products of these \emph{compressed} operators
in different replicas.  It does not identify $P_*(D_v^a)^2P_*$ with
$(B_v^a)^2$.

\begin{lemma}[Replica cross terms]\label{lem:replica-cross}
Fix $S$ and a sufficiently large finite $\beta$.  Let $\omega$ be any
translation-invariant equilibrium state of $j$ unperturbed replicas.
For distinct replicas $r,s$, axes $a,b$, and finitely supported real
slot coefficients $v,w$ with $\sum_i v_i=\sum_i w_i=0$,
\[
 \omega\bigl((B_v^a)^{[r]}(B_w^b)^{[s]}\bigr)=0.
\]
The same holds for every finite linear combination of translates of
such products.
\end{lemma}

\begin{proof}
Write $O=(B_v^a)^{[r]}(B_w^b)^{[s]}$ and choose a fixed cube of side
$d$ containing both spatial supports.  Let $\rho_L$ be the restriction
of $\omega$ to a cube of side $L$, and let $\gamma_L$ be the product
of the $j$ periodic zero-field Gibbs density matrices on that cube.
With $D(\rho\Vert\gamma)=\Tr\rho(\log\rho-\log\gamma)$,
\begin{equation}\label{eq:replica-relative-entropy}
 \frac1V D(\rho_L\Vert\gamma_L)\longrightarrow0.
\end{equation}
Indeed, expanding $\log\gamma_L$ gives $j$ times the finite-volume
pressure minus the entropy density and interaction density of
$\rho_L$.  Their limits cancel by the equilibrium variational identity;
periodic bonds at the boundary change the interaction by $o(V)$.

Fix an integer $R\ge d$ and take $L$ through even multiples of $R$.
Partition the cube into cells of side $R$, and let $\mathcal T$ be the
translations of the chosen support cube lying wholly within individual
cells.  In particular,
\[
 \frac{|\mathcal T|}{V}
       =\left(\frac{R-d+1}{R}\right)^3.
\]
Write $B_v^a=\sum_xc_xS_x^a$ and $B_w^b=\sum_xd_xS_x^b$, as in
\eqref{eq:compressed-zero-sum}.  Choose
$C_O\ge\max\{1,S\|c\|_1\|d\|_1\}$ and put
\[
 O'=\frac1{C_O}\sum_{y\in\mathcal T}\tau_y(O).
\]
All operators in this sum can be measured simultaneously, using axis
$a$ in replica $r$ and axis $b$ in replica $s$.  Conditional on any
outcome in replica $s$, denote the value of
$\tau_y(B_w^b)^{[s]}$ by $b_y$, so $|b_y|\le S\|d\|_1$.
As a function of the first replica's outcomes, $O'$ is then
$\sum_z f_zS_z^a$, where
\[
 f_z=\frac1{C_O}\sum_{y\in\mathcal T}c_{z-y}b_y.
\]
These coefficients have absolute value at most one.  Their sum in each
cell vanishes, since each translated support stays in that cell and
$\sum_xc_x=0$.  Independence of the replicas under $\gamma_L$ and
Proposition~\ref{prop:zero-field}, applied conditionally, therefore give
\begin{equation}\label{eq:replica-exponential}
 \log\Tr(\gamma_Le^{\pm\theta O'})\le C_\beta bV,
 \qquad
 b=-\log(1-R^{-5/2}),\quad\theta=\kappa b\log R.
\end{equation}

For any self-adjoint $K$, the Gibbs variational principle followed by
the Golden--Thompson inequality \cite{Golden1965,Thompson1965} gives
\[
 \Tr(\rho_LK)
 \le D(\rho_L\Vert\gamma_L)
       +\log\Tr e^{\log\gamma_L+K}
 \le D(\rho_L\Vert\gamma_L)+\log\Tr(\gamma_Le^K).
\]
Use $K=\theta O'$ or $K=-\theta O'$ according to the sign of
$\omega(O)$.  Translation invariance and
\eqref{eq:replica-exponential} yield
\[
 \theta\frac{|\mathcal T|}{C_O}|\omega(O)|
 \le D(\rho_L\Vert\gamma_L)+C_\beta bV.
\]
First let $L\to\infty$ at this fixed $R$; this order is required because
the entropy density in \eqref{eq:replica-relative-entropy} has no rate.
For large $R$, $|\mathcal T|/V$ is bounded below, and we obtain
$|\omega(O)|\le C_{\beta,O}/\log R$.  Letting $R\to\infty$ proves the
claim.  Linearity gives the final assertion.
\end{proof}

Thus, for a finite-range replica perturbation that is a sum of
single-replica terms and the cross terms in
Lemma~\ref{lem:replica-cross}, its pressure tangent depends only on
the single-replica terms.  The construction in the next section uses
this fact to study negative squares of block magnetizations.
